\documentclass[a4paper]{article}
\usepackage[pdftex]{graphicx}
\usepackage[utf8]{inputenc}
\usepackage{enumerate}
\usepackage{amssymb}
\usepackage{colortbl}
\usepackage{icomma}
\usepackage{siunitx}
\sisetup{locale=DE} 
\usepackage{geometry}
\geometry{a4paper, top=15mm, left=15mm, right=15mm, bottom=15mm,
	headsep=10mm, footskip=12mm}
\usepackage{href-ul}
\hypersetup{
	colorlinks=true,
	linkcolor=blue,
	urlcolor=blue}
\begin{document}

{\LARGE \bf $ \mathbf{F= m \cdot g} $ oder "Eine Reise zu den Sternen"}
\vspace{1,5 cm}

\noindent
{\bf Stell Dir vor: Es ist ein ganz normaler Schultag.}
Du stehst in der Pause im Schulhof, als plötzlich ein silbernes Raumschiff lautlos zwischen den Bäumen landet. Eine Rampe fährt aus, und drei freundlich lächelnde Außerirdische steigen aus.
"Komm mit uns!“, sagen sie. "Wir machen eine Forschungsreise durchs Sonnensystem – und Du darfst mitfliegen!“
Natürlich sagst Du Ja!, doch sie erklären:
"Du darfst nur sieben Dinge mitnehmen. Und bevor Du einsteigst, musst Du alles wiegen – wir wollen wissen, wie schwer Dein Gepäck auf der Erde ist.“




\begin{enumerate}[1.] 
\item Teile diese Gepäckstücke in  sinnvolle, bedingt sinnvolle und nicht sinnvolle ein:\\
\fbox{Staubsauger {\bf K}}
\noindent\fbox{Raumanzug mit Lebenserhaltungssystem {\bf J}}
\fbox{Kuscheltier {\bf Z}}
\fbox{Zahnbürste {\bf E}}\\
\fbox{Koffer voller Euros {\bf M}}
\fbox{Samentütchen mit Erdblumen {\bf  R}}	
\fbox{Notizbuch und Stift {\bf I}}\\
\fbox{Kleiner Globus als Gastgeschenk {\bf P}}	
%\fbox{Schulbuch Physik 7 {\bf I}}
\fbox{Digitalkamera {\bf T}}
\fbox{Klavier {\bf L}}
%\fbox{Schokolade {\bf T}}
%\fbox{Lieblingsbuch {\bf }}
%\fbox{Reiseapotheke {\bf }}
\fbox{Regenschirm {\bf W}}
\fbox{Kleidung {\bf U}}
\fbox{Zigarretten {\bf N}}\\
\fbox{E-Roller {\bf O}}
%\fbox{Tablet mit Offline-Sternenkarte {\bf U}}	
\fbox{Katze {\bf S}}\\

Du stellst die Gegenstände also auf eine Waage: Raumanzug, Notizbuch, Digitalkamera ...
Die Zahlen notierst Du sorgfältig – in Kilogramm. Rechne überall mit zwei gültigen Ziffern!

\item Ordne die sieben sinnvollen Sachen nach Masse. Beginne mit dem schwersten! Die Buchstaben ergeben ein Lösungswort.
\item Welche Gewichtskraft haben diese Sachen auf der Erde? Die Fallbeschleunigung auf der Erde beträgt $g_{Erde}=\SI{9,81}{\newton\per\kilogram}$

\renewcommand{\arraystretch}{2}
\begin{tabular}{|p{3 cm}|p{1,5 cm}|p{1,5 cm}|p{1,5 cm}|p{1,5 cm}|p{1,5 cm}|p{1,5 cm}|p{1,5 cm}|}
	\hline
	Gepäckstück	& & & & & & & \\
	\hline
	Masse $m$    & $\SI{80}{\kilogram}$ & & $\SI{0,75}{\kilogram}$& & $\SI{200}{\gram}$&& $\SI{5,0}{\gram}$\\
	\hline
	Gewichtskraft auf Erde $F_{g, Erde}$    & & $\SI{12}{\newton}$ & & $\SI{1,0}{\newton}$ & &  $\SI{0,20}{\newton}$ & \\
	\hline
\end{tabular}



Werte zur Kontrolle:
\fbox{$\SI{20}{\gram}$}
\fbox{$\SI{0,10}{\kilogram}$}
\fbox{$\SI{1,2}{\kilogram}$}
\fbox{$\SI{2,0}{\newton}$}
\fbox{$\SI{0,78}{\kilo\newton}$}
\fbox{$\SI{49}{\milli\newton}$}
\fbox{$\SI{7,4}{\newton}$}

\item Welche Masse haben diese Sachen auf dem Mond?

{\bf Dann geht’s los.} Die Rampe schließt sich, das Raumschiff zittert – und plötzlich spürst Du, wie Dich etwas in den Sitz drückt:
Das Raumschiff beschleunigt mit 7 g!
Du fühlst Dich siebenmal so schwer wie sonst.
Dein 40 kg-Körper "wiegt“ jetzt plötzlich so viel, als würdest Du 280 kg auf die Waage bringen. Selbst Dein Stift fühlt sich an wie aus Blei.
"Keine Sorge“, sagt der Pilot. "Das dauert nur ein paar Sekunden!“

\item Welche Geschwindigkeit hat das Raumschiff nach 100 Sekunden Beschleunigung mit 7g?

Kurz darauf wird alles ruhig – Schwerelosigkeit!
Du schwebst durch die Kabine, Dein Globus tanzt neben Dir, und dein Notizbuch schwebt durch die Luft.
Alles hat noch dieselbe Masse, aber keine Gewichtskraft mehr, weil Du und Deine Sachen gerade frei fallen – zusammen mit dem Raumschiff.

\item Der Mars sei 86,4 Millionen Kilometer von der Erde entfernt. Welche Reisegeschwindigkeit muss das Raumschiff haben, damit der Flug dorthin 24 Stunden dauert?
Gib die Geschwindigkeit bitte in $\SI{}{\kilo\meter\per\hour}$,$\SI{}{\kilo\meter\per\second}$ und $\SI{}{\meter\per\second}$ an.
\item Wie lange würde dann die Beschleunigungsphase mit 7g dauern? Gib die Zeitdauer in Sekunden, Minuten und Stunden an.

Nach einigen Stunden kommt der Mars in Sicht. Das Raumschiff bremst sanft ab und landet in einer roten Staubwolke.
Als Du aussteigst, fühlst Du Dich plötzlich leicht wie ein Ballon – denn auf dem Mars ist die Schwerkraft nur etwa ein Drittel so stark wie auf der Erde.
Ein Rucksack, der auf der Erde 9 kg wiegt, fühlt sich hier an, als hätte er nur 3 kg Gewichtskraft.


\item Auf dem Mars gilt der Ortsfaktor $g_{Mars}=\SI{3,7}{m\per s^2} .$ Wie groß ist dort die Gewichtskraft des Raumanzugs?
\item Du machst ein Selfie von Dir vor dem marsianischen Vulkan Olympus Mons und schickst es per Funk zu Deinen Eltern auf der Erde. Wie lange dauert es , bis es dort ankommt, wenn die Lichtgeschwindigkeit $c=\SI{3e8}{\meter\per\second}$ beträgt?
\item Wie würde ein Telefongespräch Mars-Erde unter diesen Bedingungen ablaufen?

Die Außerirdischen lächeln:
"Nun, Erdling – wir haben noch eine Station! Wir fliegen zu einem Exoplaneten namens Gravimurp-7 bei dem Stern Alpha Centauri. Wir möchten, dass Du dort den Ortsfaktor g bestimmst.“
Du bekommst ein Fadenpendel und Deine bekannte Ausrüstung.



\item Ihr landet auf dem Planeten "Gravimurp-7". Du merkst: Dein Pendel schwingt langsamer als auf der Erde. Du hängst Deine Digitalkamera an den Kraftmesser für die Gewichtskraft und er zeigt $F_g=\SI{1,2}{\newton}$ an.
Wie groß ist der Ortsfaktor des Exoplaneten?


Nach ein paar Messungen und Rechnungen schätzt Du:
"Hier beträgt die Fallbeschleunigung etwa $\SI{6}{\meter\per\second^2 }$– also etwas schwächer als auf der Erde, aber stärker als auf dem Mars!“
Du verstreust das Tütchen mit den Erdblumensamen auf dem feuchten Boden von Gravimurp-7.
Die Außerirdischen nicken zufrieden. 

"Sehr gut! Du hast verstanden, dass Masse immer gleich bleibt, aber die Gewichtskraft vom Ortsfaktor g abhängt.“
Dann beamen sie Dich sanft wieder in Deinen Schulhof.
Deine sieben Sachen liegen noch dort, wo Du sie hingelegt hast – nur das Staubkorn vom Mars auf Deinem Schuh erinnert Dich daran,
dass Du wirklich dort warst.

\item Zu Hause setzt Du Dich an Deinen Computer und erstellst eine Exceltabelle mit der Masse von Alltagsgegenständen  und deren Gewichtskraft auf der Erde,
während der Beschleunigung im Raumschiff, auf dem Mars und auf dem Exoplaneten Gravimurp-7:

\renewcommand{\arraystretch}{2}
\begin{tabular}{|p{1,5 cm}|p{1,5 cm}||p{3 cm}|p{2 cm}|p{2 cm}|p{2 cm}|p{2 cm}|}
	\hline
	  Ortsfaktor  in g & Ortsfaktor in  $\SI{}{\meter\per\second^2 }$& & Wasserflasche  1 kg & Katze 4 kg & Mensch 70 kg & Smartphone 0,2 kg \\
	\hline
	   1g & $\SI{9,81}{\meter\per\second^2 }$& $F_{g}$  auf Erde  & & & &\\
	\hline
	  7g & $\SI{69}{\meter\per\second^2 }$ &  $F_{g}$  in Rakete & & &  &\\
	\hline
		  0,3 g & $\SI{3,7}{\meter\per\second^2 }$& $F_{g}$  auf Mars & & & &\\
	 	\hline
	  0,6 g & $\SI{6}{\meter\per\second^2 }$ & $F_{g}$  auf Gravimurp-7	 & & & &\\
	\hline
	 
\end{tabular}
\end{enumerate} 

\fbox{
	\begin{minipage}{0.5\textwidth}
		Weitere Arbeitsblätter gibt es unter 
		
		\href{https://www.okuyakl.de}{www.okuyakl.de}
	\end{minipage}
	\hfill
	\begin{minipage}{0.4\textwidth}
		\includegraphics[width=1.5 cm]{../../viecher/zwe03}
	%	\includegraphics[width=3 cm]{qr221}
		\includegraphics[width=2 cm]{../../viecher/afanticon1}
		
\end{minipage}}

\end{document}

